\subsection{Modelo físico y simulador 6DOF}
\label{subsec:modelo-simulador}

\subsubsection*{Requisitos, convenciones e hipótesis}

El simulador debe proporcionar a ambos controladores el mismo problema físico y
la misma cadena de actuación. El estado es
$\bm{x}=(\posW,\velW,\quatWB,\omegaB)$: posición y velocidad en el marco mundo
$\frameW$, orientación del cuerpo respecto del mundo y velocidad angular en el
marco de cuerpo $\frameB$. La entrada dinámica aplicada está formada por la
fuerza total de los rotores y el momento total, ambos expresados en cuerpo. La
salida experimental incluye el estado verdadero, la observación entregada al
controlador, la referencia, el mando solicitado y la actuación aplicada.

El marco $\frameW$ sigue la convención ENU: $x$ apunta al Este, $y$ al Norte y
$z$ hacia arriba. El marco $\frameB$ es FRD: $x_B$ apunta hacia delante, $y_B$ a
la derecha y $z_B$ hacia abajo. Por ello, un rotor genera empuje en dirección
$-z_B$. La matriz $\bm{R}_{WB}(\quatWB)$ transforma vectores de cuerpo a mundo,
y $\quatWB$ se almacena en orden escalar--vector. Las magnitudes se expresan en
SI: m, s, kg, N, N\,m y rad; los ángulos solo se convierten a grados para
visualización.

% FIGURA PENDIENTE [FIG-002]: fijar los ejes ENU y FRD, la orientación del cuadricóptero y el signo del empuje

Se supone un cuerpo rígido de masa e inercia constantes, centro de masas fijo y
gravedad uniforme. Los rotores se representan mediante relaciones algebraicas y
una dinámica temporal simplificada. No se modelan aerodinámica formal,
flexibilidad ni sensores físicos. Estas hipótesis permiten aislar la comparación
de control, pero limitan cualquier extrapolación fuera del banco.

\subsubsection*{Vehículo de referencia y parámetros físicos}

La campaña v1 emplea un cuadricóptero académico de \SI{1}{kg}, gravedad
\SI{9.81}{m.s^{-2}} e inercia diagonal
$\bm{I}_B=\operatorname{diag}(0.05,0.05,0.10)\,\si{kg.m^2}$. Los cuatro rotores
ocupan las posiciones $(\pm0.17,\pm0.17,0)$~m, con sentidos de giro alternos.
Cada rotor utiliza $k_f=10^{-4}$, $k_m=10^{-6}$ y
$\omega_{\max}=\SI{1500}{rad.s^{-1}}$. El valor máximo teórico por rotor es,
por tanto, $k_f\omega_{\max}^2=\SI{225}{N}$; este límite pertenece al modelo y
no debe interpretarse como especificación de un vehículo comercial.

% CITA PENDIENTE [CIT-002]: criterio físico o fuente de referencia para masa, geometría, inercia y coeficientes de propulsión del vehículo académico
Los valores se mantienen constantes para que las diferencias observadas
procedan del controlador y del escenario, no de vehículos distintos. Sin una
identificación o una fuente de plataforma concreta, constituyen parámetros de
un caso académico controlado. Su validez se limita a comprobar coherencia
interna y reproducibilidad; no permiten afirmar representatividad cuantitativa
de un dron real.

\subsubsection*{Cinemática y dinámica del cuerpo rígido}

Sean $\bm{F}_B$ y $\bm{\tau}_B$ la fuerza y el momento aplicados en cuerpo. La
dinámica traslacional implementada es
\begin{align}
 \dot{\posW} &= \velW,\\
 m\dot{\velW} &= \bm{R}_{WB}\bm{F}_B+\bm{F}_{d,W}
                 +[0\;0\;-mg]^\mathsf{T},
 \label{eq:dinamica-traslacional}
\end{align}
donde $\bm{F}_{d,W}$ es el drag descrito en el
\apartadoref{subsec:modelo-simulador}. La rotación sigue
\begin{align}
 \dot{\quatWB} &= \tfrac{1}{2}\quatWB\otimes[0\;\omegaB^\mathsf{T}]^\mathsf{T},\\
 \dot{\omegaB} &= \bm{I}_B^{-1}
 \left(\bm{\tau}_B-\omegaB\times\bm{I}_B\omegaB\right).
 \label{eq:dinamica-rotacional}
\end{align}
La normalización del cuaternión se aplica en las evaluaciones intermedias y al
final de cada paso para contener la deriva numérica de su norma.

El estado avanza mediante Runge--Kutta explícito de cuarto orden (RK4) y paso
fijo de \SI{0.01}{s}. El mando aplicado se mantiene durante el paso. Esta
elección ofrece una discretización determinista y un coste predecible para la
campaña.
% CITA PENDIENTE [CIT-003]: propiedades y condiciones de uso de RK4 de paso fijo en integración de ecuaciones diferenciales ordinarias
No se ha realizado todavía un estudio de convergencia temporal que demuestre que
\SI{0.01}{s} sea suficiente para todas las condiciones; por ello el paso se
declara parte del protocolo y no una garantía general de exactitud.

\subsubsection*{Actuadores, mezclador y límites de actuación}

Para el rotor $i$, la velocidad aplicada determina empuje y par de reacción
mediante
\begin{equation}
 T_i=k_f\omega_i^2,\qquad Q_i=s_i k_m\omega_i^2,
 \label{eq:rotor-cuadratico}
\end{equation}
donde $s_i\in\{-1,1\}$ representa el sentido de giro. La fuerza es
$\bm{F}_{i,B}=[0\;0\;-T_i]^\mathsf{T}$ y su contribución al momento incluye
$\bm{r}_{i,B}\times\bm{F}_{i,B}$ y $[0\;0\;Q_i]^\mathsf{T}$.
% CITA PENDIENTE [CIT-004]: fundamento de la ley cuadrática empuje/par de rotor y asignación de fuerzas en un cuadricóptero

El mezclador construye una matriz de asignación con estas contribuciones y
resuelve la demanda de empuje colectivo y momentos. Si la solución exige
empujes negativos, desplaza el conjunto para conservar la distribución relativa;
después limita cada rotor al intervalo realizable y marca la degradación del
mando. La conversión final usa $\omega_i=\sqrt{T_i/k_f}$.

La dinámica de cada rotor combina un retardo puro discretizado y un lag de
primer orden. En el perfil nominal, la constante de tiempo es \SI{0.03}{s} y el
retardo \SI{0.01}{s}; el perfil combinado aumenta ambos a \SI{0.05}{s} y
\SI{0.02}{s}. La saturación se aplica a $[0,\omega_{\max}]$ y queda registrada.
% CITA PENDIENTE [CIT-005]: justificación del modelo de primer orden, el retardo y sus valores para los actuadores simplificados
Estos parámetros son elecciones operativas de v1, no una identificación de
motor, hélice y variador concretos.

\subsubsection*{Perturbaciones, observación y seguridad}

El drag lineal se calcula en cuerpo a partir de la velocidad relativa al viento,
$\bm{v}_{r,B}=\bm{R}_{WB}^{\mathsf{T}}(\velW-\bm{v}_{w,W})$, y vuelve a mundo:
\begin{equation}
 \bm{F}_{d,W}=\bm{R}_{WB}\left(-\bm{D}_B\bm{v}_{r,B}\right).
 \label{eq:drag-lineal}
\end{equation}
Su signo es disipativo. Los perfiles emplean diagonales base
$(0.10,0.10,0.05)$ o elevadas $(0.20,0.20,0.10)$, viento constante y ruido
gaussiano independiente sobre posición y velocidad. El controlador recibe la
observación perturbada; la dinámica y la evaluación conservan el estado
verdadero.

% CITA PENDIENTE [CIT-006]: alcance y limitaciones de representar resistencia, viento y ruido mediante modelos lineales, constantes y gaussianos
Estos mecanismos permiten variar de forma controlada disipación, sesgo externo y
calidad de observación, pero no reproducen turbulencia ni sensores reales. Los
perfiles concretos son niveles heurísticos del banco y se someten al mismo
tratamiento para todos los controladores.

El episodio termina si $z_W<0$, si roll o pitch superan \SI{1.256}{rad}, si la
saturación persiste más de \SI{2}{s}, si aparecen valores no finitos o si alguna
componente de posición o velocidad excede los límites internos de \SI{100}{m} y
\SI{50}{m.s^{-1}}. Los \SI{1.256}{rad} separan una actitud extrema de la región
de operación prevista, pero no equivalen a un límite certificado de vuelco.
% CITA PENDIENTE [CIT-007]: criterio defendible para el límite de actitud y los umbrales de seguridad de la campaña

\subsubsection*{Trayectorias y escenarios reproducibles}

Las cuatro familias de entrenamiento ejercitan capacidades diferentes.
\texttt{hold} exige mantener posición y yaw; el círculo combina aceleración
centrípeta y avance continuo; Lissajous introduce movimiento periódico acoplado
en varios ejes; y \texttt{waypoint} exige desplazamientos finitos, frenado,
asentamiento y permanencia antes de avanzar. Las trayectorias analíticas entregan
posición, velocidad y aceleración suaves. Las misiones por puntos usan perfiles
trapezoidales o triangulares con parada intermedia y no escalones crudos de
posición.

Cada YAML fija vehículo, estado inicial, trayectoria, controlador,
perturbaciones, semilla, periodos, terminación y directorio de salida. El estado
inicial coincide con la referencia en $t=0$, con velocidades nulas y actitud
nivelada, de modo que el conjunto principal evalúa seguimiento y no captura
desde un error inicial arbitrario. La razón experimental de cada geometría y
conteo se desarrolla en el \apartadoref{subsec:datos-metricas}; las
composiciones y lemniscatas se reservan para OOD.

\subsubsection*{Flujo multirrate de simulación}

La física se actualiza cada \SI{0.01}{s}, el control cada \SI{0.02}{s} y la
telemetría cada \SI{0.1}{s}. En cada instante se obtiene la referencia, se forma
la observación y, cuando corresponde, se calcula un mando nuevo. El mezclador y
los actuadores producen las fuerzas aplicadas; RK4 avanza el estado y la
telemetría captura las variables en su propio calendario. Entre dos ciclos de
control se aplica retención de orden cero del mando solicitado.

% FIGURA PENDIENTE [FIG-003]: representar el orden de referencia, observación, control, mezclador, actuadores, RK4 y telemetría con sus tres periodos
% CITA PENDIENTE [CIT-008]: justificación de la retención de orden cero y de la separación multirrate en simulación digital de control
La separación evita recalcular la política a la frecuencia de integración y
reduce el volumen de telemetría sin cambiar la dinámica. Los valores concretos
son parte del diseño v1; falta todavía un análisis de sensibilidad frente a
otros periodos.

\subsubsection*{Telemetría, verificación y trazabilidad}

Cada muestra conserva tiempo, estado verdadero, observación, referencia, mando
solicitado, objetivo de rotor, actuación aplicada y causa de terminación. De este
modo pueden distinguirse error de seguimiento, ruido observado, degradación del
mezclador, retardo y saturación. El resumen añade estadísticos con unidades y
metadatos: escenario, semilla, controlador, configuración original y efectiva,
comando, versión del paquete, plataforma, estado Git y hashes del escenario y
de \texttt{uv.lock}.

Las pruebas automáticas comprueban, entre otros contratos, norma del cuaternión,
signo del empuje ENU/FRD, dinámica, RK4, mezclador, saturación, perturbaciones,
ejecución multirrate, terminaciones y telemetría. Estas pruebas verifican que la
implementación sea coherente con sus ecuaciones y contratos. No validan la
fidelidad del modelo frente a datos de vuelo; tal validación requeriría
identificación y comparación experimental externa.

\subsubsection*{Elección del entorno de implementación}

Python permite expresar el modelo con operaciones vectoriales, declarar
escenarios legibles y compartir el mismo entorno con el entrenamiento en
PyTorch. \texttt{uv} fija dependencias y comandos, mientras YAML separa los
parámetros experimentales del código. La elección facilita inspección,
automatización y trazabilidad, requisitos centrales de este banco.

No se afirma que Python sea la opción de mayor rendimiento ni que el simulador
sea apto para tiempo real. La ejecución interpretada y el propósito académico
acotado limitan la escala y fidelidad alcanzables. La comparación con otras
plataformas corresponde al estado del arte; la organización interna completa
se remite al repositorio\footnote{\url{https://github.com/ChristianGonper/modelado-control-dron-v2}}.
